\chapter{PEFT, LoRA, adapter, MoE: 얼마나 작은 상태를 바꿀 것인가}
\label{bg:chapter-peft}

\section{Trainable state를 줄이는 이유}

Full fine-tuning은 모든 weight에 gradient와 optimizer state를 만든다. 사용자·도메인별
variant가 많아지면 compute뿐 아니라 checkpoint storage와 model residency가 폭증한다.
\concept{peft}{PEFT}는 base model을 고정하고 작은 subset 또는 추가 module만 학습해 이
비용을 줄인다. 하지만 trainable parameter 비율과 total system cost를 구분해야 한다.
Base weight는 여전히 forward/backward에 참여하며 activation과 input data를 읽는다.

\concept{trainable-state}{학습 가능 상태}는 gradient를 받아 갱신되는 상태다. 반면
\concept{serving-artifact}{서빙 산출물}은 실제 wake inference에 배포되는 base checkpoint,
delta, routing table, index의 조합이다. Sleep job이 작은 delta만 만들더라도 수백만 사용자의
delta를 catalog, cache, prefetch, revoke하는 serving system은 클 수 있다.

\section{LoRA의 최소 수식}

\concept{lora}{LoRA}는 고정 weight $W_0\in\mathbb{R}^{d_{out}\times d_{in}}$에 low-rank
delta를 더한다 \parencite{hu2021lora}.

\begin{equation}
  W = W_0 + \Delta W,
  \qquad \Delta W = \frac{\alpha}{r}BA,
  \label{eq:bg-lora}
\end{equation}

$A\in\mathbb{R}^{r\times d_{in}}$, $B\in\mathbb{R}^{d_{out}\times r}$이고
$r\ll \min(d_{in},d_{out})$다. full matrix의 $d_{out}d_{in}$개 대신
$r(d_{in}+d_{out})$개를 학습한다. 예를 들어 $d_{in}=d_{out}=4096$, $r=16$이면 한
projection의 trainable parameter는 약 16.8M에서 131K로 줄어 약 128배 작다.

그러나 다음 비용은 남는다.

\begin{itemize}
  \item base layer를 통과하는 forward와 backward activation,
  \item LoRA A/B gradient와 optimizer state,
  \item 여러 target module(Q/K/V/O, MLP)에 대한 delta,
  \item user/domain별 adapter 저장·인덱싱·authentication,
  \item base checkpoint가 바뀔 때 compatibility 재검증,
  \item 여러 adapter를 합성할 때 발생하는 conflict.
\end{itemize}

\section{Adapter module과 routing}

\concept{adapter}{어댑터}는 layer 사이에 작은 bottleneck module을 추가한다. 전형적인 형태는

\begin{equation}
  h' = h + W_{up}\,\sigma(W_{down}h)
\end{equation}

이다. Base를 고정할 수 있고 adapter별 version이 명확하다. 대신 inference graph에 추가
operator가 생기며, 어떤 adapter를 언제 load할지 routing해야 한다. adapter가 HBM에 없으면
storage-to-HBM load latency가 first-token latency에 나타날 수 있다. 여러 작은 adapter가
random access workload를 만들면 총 byte보다 IOPS와 metadata lookup이 병목일 수 있다.

\section{Mixture of Experts는 PEFT와 다른 축이다}

\concept{mixture-of-experts}{Mixture of Experts(MoE)}는 token마다 router가 일부 expert만
활성화하는 sparse architecture다. 전체 parameter capacity는 크지만 active FLOPs는 일부다.
MoE가 자동으로 개인 memory system이 되는 것은 아니다. Expert가 pretrained/fixed인지,
sleep 중 새 expert를 추가하는지, user-specific routing을 학습하는지에 따라 의미가 달라진다.

Sleep-time 관점에서 MoE는 세 기회를 준다. (1) stable domain knowledge를 expert로 격리,
(2) base 전체가 아닌 selected expert만 update, (3) cold expert를 lower memory tier에 offload.
동시에 all-to-all communication, expert imbalance, cache miss, expert proliferation이라는 시스템
문제를 만든다.

\begin{table}[htbp]
\centering
\caption{Full FT, LoRA, adapter, external memory의 비교}
\label{tab:bg-peft-comparison}
\small
\begin{tabularx}{\textwidth}{L{0.15\textwidth}YYYYY}
\toprule
방법 & trainable state & optimizer state & wake access & rollback & 핵심 병목 \\
\midrule
Full FT & 전체 weight & 전체 규모 & 직접 weight & checkpoint 단위 & HBM/compute/interference \\
LoRA & low-rank delta & delta 규모 & base+delta matmul & delta version & adapter catalog/cache \\
Adapter & 작은 module & module 규모 & extra module & module version & routing/load latency \\
MoE update & selected experts & expert subset & sparse routing & expert/version & all-to-all, imbalance \\
External memory & text/vector/graph & 선택적 & retrieval+context & record/version & index, context pollution \\
\bottomrule
\end{tabularx}
\end{table}

\section{Selective promotion이라는 hybrid}

모든 episode를 weight로 넣지 않고 external system-of-record에 먼저 보관한다. 충분한 future
reuse, stability, cross-user generality, safety evidence가 쌓인 item만 adapter나 periodic global
refresh로 승격한다. 이 구조는 raw episode의 provenance와 delete를 유지하면서 high-reuse
knowledge의 repeated retrieval cost를 줄일 수 있다.

Promotion break-even의 최소 형태는 다음과 같다.

\begin{equation}
  C_{\mathrm{compile}}+C_{\mathrm{validate}}+C_{\mathrm{store}}
  < N_{\mathrm{future}}\left(C_{\mathrm{retrieve}}-C_{\mathrm{parametric\ access}}\right).
  \label{eq:bg-promotion-break-even}
\end{equation}

$N_{\mathrm{future}}$가 작거나 base model churn이 빠르면 승격은 손해다. 따라서 adapter는
``항상 더 효율적인 기억''이 아니라 later reuse가 충분할 때만 이기는 compiled cache에
가깝다.

\begin{cautionbox}[작은 delta의 함정]
LoRA file이 작다는 사실만으로 user-specific parametric memory가 scale한다고 결론 내릴 수
없다. adapter 수, active working set, load pattern, validation cost, base-version migration,
deletion SLA를 함께 측정해야 한다.
\end{cautionbox}
